%Copyright 2014 Jean-Philippe Eisenbarth
%This program is free software: you can 
%redistribute it and/or modify it under the terms of the GNU General Public 
%License as published by the Free Software Foundation, either version 3 of the 
%License, or (at your option) any later version.
%This program is distributed in the hope that it will be useful,but WITHOUT ANY 
%WARRANTY; without even the implied warranty of MERCHANTABILITY or FITNESS FOR A 
%PARTICULAR PURPOSE. See the GNU General Public License for more details.
%You should have received a copy of the GNU General Public License along with 
%this program.  If not, see <http://www.gnu.org/licenses/>.

%Based on the code of Yiannis Lazarides
%http://tex.stackexchange.com/questions/42602/software-requirements-specification-with-latex
%http://tex.stackexchange.com/users/963/yiannis-lazarides
%Also based on the template of Karl E. Wiegers
%http://www.se.rit.edu/~emad/teaching/slides/srs_template_sep14.pdf
%http://karlwiegers.com
\documentclass[16pt]{scrreprt}
\usepackage{listings}
\usepackage{float}
\usepackage{setspace}
\usepackage{url}
\usepackage{longtable}
\usepackage{booktabs}
\usepackage{underscore}
\usepackage{graphicx, subfig}
\usepackage[bookmarks=true]{hyperref}
\usepackage[utf8]{inputenc}
\usepackage[english]{babel}

\hypersetup{
    bookmarks=false,    % show bookmarks bar?
    pdftitle={Software Requirement Specification},    % title
    pdfauthor={Jean-Philippe Eisenbarth},                     % author
    pdfsubject={TeX and LaTeX},                        % subject of the document
    pdfkeywords={TeX, LaTeX, graphics, images}, % list of keywords
    colorlinks=true,       % false: boxed links; true: colored links
    linkcolor=blue,       % color of internal links
    citecolor=black,       % color of links to bibliography
    filecolor=black,        % color of file links
    urlcolor=purple,        % color of external links
    linktoc=page            % only page is linked
}%
\def\myversion{1.0 }
\date{}
%\title{%

%}
\usepackage{hyperref}
\DeclareOldFontCommand{\bf}{\normalfont\bfseries}{\mathbf}
\begin{document}

\begin{flushright}
    \rule{16cm}{5pt}\vskip1cm
    \begin{bfseries}
        \LARGE{SOFTWARE REQUIREMENTS\\ SPECIFICATION}\\
        \vspace{0.8cm}
        for\\
        \vspace{0.8cm}
        iSport\\
        \begin{figure}[h]
		\flushright
  		\includegraphics[width=.2\textwidth]{logo.png}
		\end{figure}
        \vspace{0.8cm}
        \normalsize{Version \myversion approved}\\
        \vspace{1.0cm}
        Prepared by\vspace{0.5cm} \\ 
        \normalsize{1650940 Jiang Xiaohu\\1650932 Xu Jingnan\\1651058 Wang Yicheng}
        \\
        \vspace{1.1cm}
        \normalsize{School of Software Engineering\\ Tongji University}\\
        \vspace{1.0cm}
        \today\\
    \end{bfseries}
\end{flushright}

\tableofcontents


\chapter*{Revision History}

\begin{center}
    \begin{tabular}{|c|c|c|c|}
        \hline
	    Name & Date & Reason For Changes & Version\\
        \hline
	    Jiang Xiaohu & 2019.11.1 & Finish Introduction Part  & v1.0\\
        \hline
	    31 & 32 & 33 & 34\\
        \hline
    \end{tabular}
\end{center}

\chapter{Introduction}

\section{SRS Purpose}
The purpose of this document is to present a detailed description of iSport. It will explain the purpose and features of iSport, the interfaces of the iSport, the functional and nonfunctional requirements of iSport, what iSport will do, and the constraints under which it must operate and how the system will react to external stimuli. This document is intended for both the stakeholders and the developers of the system and will be proposed to the clients for its approval.
%This SRS (Software Requirements Specification) document aims to identify the product--iSport whose software requirements are specified in this 
%document, including the revision or release number. And the scope of iSport is covered by this SRS, too, particularly if this SRS describes only 
%part of the system or a single subsystem.$>$

\section{Product Scope}
This software system will be a web based system for sports fans, professional athletes, and patients who are under recovering training. \\
\\
This system will be designed to maximize the exercising efficiency by providing tools to assist in checking and correcting user's wrong postures and recommending training courses customized for users, which would otherwise have to be expensive, time-consuming and labor intensive. By maximizing the user’s training efficiency and convenience the system will meet the needs of sports fans, athletes and injured patients while remaining easy to understand and use.\\
\\
More specifically, this system is designed to allow a user to imitate the standard exercising postures while observe and correct their mistakes simultaneously with the help of a website. \\
\\
The software will collect some professional courses in the database, including static and dynamic trainings which means doing exercise according to a set of images or a video and iSport will recommend suitable trainings for users on the basis of their training performance. Courses are classified into exercising courses and recovering courses, aiming to help athletes and patients respectively.\\
\\
Both visual and audio notification are used in every course of the system to provide eye-catching, user-friendly and clear instructions; the feedback of one's training is proposed once the training is over and the report can be browsed in the report page.\\
\\
The selection and deletion of one user's favorable course is supported in personal information webpage and one can comment training he/she has taken on comment webpage to provide suggestions to other users.\\
\\
 The personal information registering, changes is allowed via the application options. The system also contains a relational database containing a list of users, training images and videos.

\section{Intended Audience}
This Software Requirements document is intended for:
\begin{itemize}
	\item Developers who can review project’s capabilities and more easily understand where their efforts should be targeted to improve or add more features to it (design and code the application – it sets the guidelines for future development).
	\item Project testers can use this document as a base for their testing strategy as some bugs are easier to find using a requirements document. This way testing becomes more methodically organized.
   \item End users of this application who wish to read about what this project can do.
   \item Clients who delegate the software development to our team, and can check if their requirements are perfectly understood by the developers and whether the functional and nonfunctional requirements are entirely meet. And they can modify some requirements according to this document in later stage. 
   \item Project managers who translate the clients' requirement to the programmers and supervise them to implement the requirements mentioned in the documents and if the clients modify the requirements, the project managers have to negotiate the changes with the clients and modify the document.
   \item Marketing staff who are responsible to sale and prompt the software should be clear about what iSport can do, what iSport's advantages are, what iSport's competitive power is.
   \item Service staff who are responsible to solve the customer's problems, they have to know what iSport can do and feedback the technical problems to developers.
   \item Document writers who are responsible for writing the rest documents in later development stages should follow the requirements specified in this SRS.
\end{itemize}
The rest part of this SRS document contains the overall description of iSport, and iSport's specific, nonfunctional and other requirements which are shown in chapter\ref{Overall Description}, chapter\ref{Specific Requirements}, chapter\ref{Other Nonfunctional Requirements} and chapter\ref{Other Requirements} respectively.
We suggest the readers begin with the overview sections and proceeding through the sections that are most pertinent to them:
\begin{itemize}
	\item Developers and project testers are recommended to focus on specific and nonfunctional requirements part because theses parts will lead them to build qualified, safe and satisfying application and theses parts are all related to coding (construction and verification stage).
	\item Clients, project managers and document writers should focus on the entire document since they are responsible for all the requirements specified in this paper.
	\item Marketing staff have to focus on the functional part of iSport.
\end{itemize}


\section{Definitions, Acronyms, and Abbreviations}
\begin{longtable}{|p{1.9in}|p{4in}|c|}
xxxxx & xxxxxx  \kill
\caption{Definitions\label{simple}}\\ \hline
\multicolumn{3}{|c|}{\bf Definitions, Acronyms, and Abbreviations}\\ \hline
\endfirsthead
\caption[]{(continued)}\\ \hline
\multicolumn{3}{|c|}{\bf Definitions, Acronyms, and Abbreviations (continued)}\\
\hline
\endhead
\hline
\multicolumn{3}{|c|}{\bf Continued $\ldots$}\\
\hline
\endfoot
\hline
\multicolumn{3}{|c|}{\bf The End}\\
\hline
\endlastfoot
Term & Definitions  \\
\hline
User & Someone who interacts with iSport including sports fans, athletes and  injured patients who need recovery training.\\  \hline  
Sports fans & One of iSport's potential customers who love sports and want to get professional instructions when exercising. Some of them may can't afford the expense of personal coaching or don't have time to go to the gym. \\ \hline
Athletes & One of iSport's potential customers who want to get real-time exercising feedback to improve their performance or who want to get some relaxing training in their spare time to keep a good competitive state.\\  \hline
Injured Patients & One of iSport's potential customers who need recovering training after some treatments, e.g. surgeries. On the one hand, some of them may can't afford the doctor's expensive medical instructions for recovering training. on the other hand, there is no enough doctors or nurses who can instruct and supervise the patients' recovering exercising. But without professional instructions training can be useless or even leads to secondary trauma.\\  \hline
Admin/Administrator & System administrator who is given specific permission for managing and controlling the system, e.g. updating the user's information, uploading new training courses.\\ \hline
User Info & User's basic information including user's avatar, account name, tel-number and email address.\\ \hline
Courses & Training courses including normal exercising training and recovering training.\\ \hline
Normal Courses & Training courses which serve the sports fans and athletes.\\ \hline
Recovering Courses & Training courses which serve the injured patients.\\ \hline
Static Courses & Training courses which instruct the users photo by photo.\\ \hline
Dynamic Courses & Training courses which instruct the users according to a standard video.\\ \hline
Appraisal Subsystem & Remark the user's performance by using a grade from 0 -100\\ \hline
Comment Subsystem & User comment on the training courses they have taken to provide reference for other users.\\ 
\hline 
Recommendation Subsystem & A subsystem which will provide some courses for users according to their recent performance.\\ 
\hline
Exercise Tips & There will be sports tips in the webpage of iSport to prevent users from athletic injuries.\\ 
\hline 
Sport Report & A web page to feedback the user's exercising performance.\\ 
\hline 
Audio Notification & An audio notification will be shown when the user is doing exercise to encourage the user to hold on or notify the user to correct their postures.\\ 
\hline 
Visual Notification & A visual notification will be shown when the user is doing exercise, if the user's posture is standard, then the web-frame will turn green to suggest the user to hold on, otherwise the web-frame will be red.\\ 
\hline 
DataBase & A relational database containing a list of user info, training images and videos.\\ 
\hline 
Detection Subsystem & Subsystem to detect the user's postures and draw the user's skeleton. The main model of detection subsystem is PoseNet.\\ 
\hline 
Comparison Subsystem & Subsystem to compare the postures of the user and that of the standard. The subsystem aims to check if the user pass the posture.\\ 
\hline 
Correction Subsystem & Subsystem to calculate where the postures' wrong part are, e.g. left-arm, right-leg, head.\\ 
\hline 
Clients & Group who delegate the development of iSport to the developers and will take charge of the later management of iSport.\\ \hline
Developers & Develop team including project managers, programmers, testers who are responsible for the development of iSport and its later mainteinance and updating.\\ 
\hline 
\end{longtable}

\section{Document Conventions}
This document follows MLA Format[1]. Bold-faced text has been used to emphasize sectionand sub-section headings. Highlighting is to point out the references of tables and figures. And italicized text is used to label and recognize special characters or terminologies\\

The SRS paper is written by latex[4], the packages we used to format the document including \small{$longtable, graphicx, subfig, utf8-inputenc,hyperref$ }and so on.\\
And we modified the IEEE SRS latex template from Jean-Philippe Eisenbarth's github[5] according to the IEEE standard[2] mentioned in Frank F Tusi'book[3].


\section{References and Acknowledgements}
\subsubsection{Standard Reference}
The standard we have followed are as follows:\\

$[1]$ T. Russell, A. Brizee, E. Angeli, and R. Keck, “Mla formatting and style guide,” The Purdue OWL, 2010.\\

$[2]$ I. S. E. S. Committee et al., “Ieee recommended practice for software re- quirements specifications,” IEEE organization, 1998.\\


$[3]$ F. F. Tsui, O. Karam, and B. Bernal, Essentials of software engineering. Jones Bartlett Learning, 2016.

\subsubsection{Writing Tools Reference}
The writing tools we have used are as follows:\\

$[4]$L. Lamport, LATEX: a document preparation system: user’s guide and reference manual. Addison-wesley, 1994.\\

$[5]$ J.-P. Eisenbarth, “Srs latex template under ieee standard,” \\http://https://github.com/jpeisenbarth/SRS-Tex.


\chapter{Overall Description}
\label{Overall Description}

\section{Product Perspective}

ISport is a web system developed by isport team, aimming at posture correcting with the assist of camera.\\

The context diagram in the folowing figure illustrates the external entities and system interfaces. The system is expected to evolve over at least three releases, ultimately allowing for complete streamlining of the posture correcting process, fitness classes and rehabilitation classes for learning.

\section{Product Functionality}


ISport contains these following key features:\\

Let the user register their account of the website through their mobile phone or e-mail.\\

Let the user log in to the system by correctly complete the log-in form. If the username is not in the database, the user will be prompted to register his/her account.\\

User can modify his/her account information after logging in, including the username, e-mail address, phone number, gender, age and so on.\\

Show all the sorted fitness courses in a list for user to choose in the fitness courses view page. \\

Show a list of courses that the user has selected to learn before\\

Show a list of recommended courses for the user by user's previous behavior \\

Show the comment and evaluation of each courses in their detailed page\\

Provide the add button for user to add a course he/she want to learn in the future and add the course to his/her courses list.\\

Provide the delete button for user to remove a course he/she doesn't want to learn and remove the course from his/her courses list.\\


Generate user's exercise report for watching and analyzing in the report page\\

Give the user a space to leave his/her evaluation about the courses learnt before\\

Capture and Collect the posture data of the user in front of the computer and store the data in the server.\\


Compare user's posture with the standard one and calculate the similarity between the postures.\\


Show tips on the user's screen to notice the wrong posture of the user and help to correct them.\\


Rate the user's posture on a scale of zero to ten to let the user know whether his posture is standard or not\\


System administrator can upload new courses to the system for users to choose and learn.\\


\section{Users and Characteristics}

\begin{center}
    \begin{tabular}{p{5cm}p{11cm}}
        \hline
	    Users & Desc\\
        \hline
	    People who want to correct their posture &  Normal user is expected to register and log in the system, upload their posture and get the feedback advice from the system\\
        \hline
	    People who want to be fitness & People who want to be fitness is expected to register and log in the system, choose fitness classes for themselves and learn it on the website\\
        \hline
        People who need rehabilitation & People who need rehabilitation is expected to register and log in the system, choose fitness classes for themselves and learn it on the website\\
        \hline
        System Administrator & System Administrator has the privilege to update posture information in the database. The Administrator does not directly interact with the website\\
        \hline
    \end{tabular}
\end{center}

\section{Operating Environment}

\subsubsection{Hardware}


\noindent1. Server\\

CPU: Intel Dual-Core 2.4GHz\\

Memory: 8GB\\

External Storage: 500G SSD\\

Quantity: 1\\

\noindent 2. Client(Minimum Configuration)\\

CPU: Single-Core 1GHz\\

Memory: 2G\\

External Storage: 20GB HDD\\


\subsubsection{Software}


\noindent 1. Server\\

Operating System: Ubuntu 18.04 LTS\\

Database: MySQL\\

Software and Library : Python, Spring, TensorFlow, Nginx\\


\noindent 2. Client\\

Operating System: Windows/macOS/Linux/Android/IOS\\

Software: Browser\\


\begin{center}
    \begin{tabular}{p{7cm}p{7cm}}
        \hline
	    Recommand Browser & Version\\
        \hline
	    Google Chrome &  44+\\
        \hline
	    Mozilla Firefox & 40+\\
        \hline
        Apple Safari & 7+\\
        \hline
        Microsoft Edge & 12+\\
        \hline
        Microsoft Internet Explorer & 11+\\
        \hline
        Opera Opera & 31+\\
        \hline

    \end{tabular}
\end{center}

\section{Design and Implementation Constraints}

\noindent 1. Memory: Server will have 500GB internal hard drive. Softwares and database cannot exceed this amount. System administrator must notice this limitation. And each user should follow the rule that the video data uploaded each time shall not be greater than 100M\\


\noindent 2. Language requirements: software must be multilingual, including the following languages: English, Chinese\\


\noindent 3. Number of user: each video uploaded must be one person. Each time this system can only deal with one person's data.\\

\section{User Documentation}

Along with this system: iSport, a user manual need to be written to help users understand how to operate the system. It would be written for nontechnical individuals and the level of content would differ considerably from a system administration guide, which is more detailed and complex. The user manual would follow common user documentation styles to be simple.\\

Trying to use step-by-step instructions for users who firstly log in to the website, by showing messaging structures, quick references, tips and glossary of terms.\\

User document can be written in HyperText Markup Language (HTML) or Portable Document Format (PDF) , which must describe the use of the software system.

\section{Assumptions and Dependencies}

It is assumed that the website will work correctly with every third-party operating
system and compatible across all of the major browsers.\\

Assumed that the web server always runs well without down or not responding.\\


ISport provides two kinds of classes for user to choose: fitness classes and rehabilitation classes.\\

Website visitors who have not been registered are only allowed to watch some demo videos, they are not allowed to upload their posture video data before registering.\\

In addition to test their posture and attend into courses, members are also allowed to update their information, but they have to log in first.\\


\chapter{Specific Requirements}
\label{Specific Requirements}
\section{External Interfaces Requirements}
\subsection{User Interfaces}
\begin{itemize}
    \item GUI Standard: Material Design is a design language that Google developed in 2014. Expanding on the "card" motifs that debuted in Google Now, Material Design uses more grid-based layouts, responsive animations and transitions, padding, and depth effects such as lighting and shadows.
    \item font Standard: Google Sans Font, agree with Material Design
    \item button Standard: Material Design
    \item color scheme: Dark, Light, changes in different system color scheme. Nowadays, every OS has at least 2 types color scheme, someone may prefer dark mode to be less distract.
    \item screen layout/resolution: adapt to user's device
    \item function, navigation links: in login page, user can go to homepage or register page by router. After entering homepage, user can see a navigation bar, in which user can jump to every sub-page without difficulty.
    \item message display: every dangerous change, such as changing password, username, phone number, the web page will show a message notification, on the righttop. If user does some wrong gestures, it will show a message notification and bell a sound to warn.
\end{itemize}
\subsection{Hardware Interfaces}
\begin{itemize}
    \item supported device:\\
        \begin{itemize}
            \item mobile phones, including Android and iOS devices, with camera authorization.
            \item computers with camera or laptop, with camera authorization.
        \end{itemize}

    \item deployed device:\\
        A cloud web server in Ali Cloud, with 4 gigabytes RAM and dual-core CPU and 1TB disk.\\
        We deployed our web service, including database, backend on one server, the user amount now is not a large number, in some day if the user group grows larger and larger, we may use multiple servers and consider disturbed system.  
\end{itemize}
    Our product recall the interface between camera, to update data which contains the gestures user does to our server to analysis whether what the user does corresponds to criterion.
\subsection{Software Interfaces}
\begin{itemize}
    \item Third part APIs we use
        \begin{itemize}
            \item OS: our service requires the OS grant the permission to recall camera of user divices.
            \item database: the database of our service is generally MySQL. Java Database Connectivity (JDBC) is an application programming interface (API) for the programming language Java, which defines how a client may access a database. It is a Java-based data access technology used for Java database connectivity. It is part of the Java Standard Edition platform, from Oracle Corporation. It provides methods to query and update data in a database, and is oriented towards relational databases. A JDBC-to-ODBC bridge enables connections to any ODBC-accessible data source in the Java virtual machine (JVM) host environment. To make jwt works correctly, we use a key-value dababase called redis, storing outdated token and blocking them to avoid a situation that a user just changed his password, but his token is still not outdated, because jwt cannot forces a token to be outdated.
            \item Springboot: the backend programm language is Java, our backend is based on Springboot and developed in form of microservice architecture. Microservice architecture is a development trend and arranges an application as a collection of loosely coupled services. In a microservices architecture, services are fine-grained and the protocols are lightweight. With APIs provided by Springboot, the development goes efficiently and standard.
            \item PoseNet: our identification technologies is provided by a open-source project named PoseNet, PoseNet contains a standalone model called PoseNet, as well as some demos, for running real-time pose estimation in the browser using TensorFlow.js.
        \end{itemize}
    \item APIs we provide
        \begin{itemize}
            \item user register API\\
                This API provides the function that a new user becomes a new menber of iSport.\\
                API checks:
                \begin{itemize}
                    \item whether the username has registered, 
                    \item whether the phone or email is empty
                \end{itemize}
                if pass all restrict, the new user owns his data in our database.
            \item user login API\\
                This API provides the function that a user launch a new session.\\
                API checks:
                \begin{itemize}
                    \item whether this user has register.
                    \item the correction of password
                \end{itemize}
                if pass all restrict, the user now has logged in. He can begin to exercise.
            \item all courses API\\
                This API recieves a integer send by client, requests backend to return all courses with the type user type in. There two types of courses in iSport, one is rehabilitation course and the other is training course.
                API checks the type typed in by user, an integer, to return. Only 0 and 1 is valid, 1 is for rehabilitation, and 0 is for training.
            \item course content API\\ 
                This API recieves an ID of a course, and return all photography of this course, the provide a basic message to users.\\
                API checks the input, make sure that the input is not an empty input, then query in database, if the ID is invalid(e.g. not a course with this ID), the API throw an exception. If not, API respond with all photography about this course in database.
            \item get user name\\
                This API revives the ID of user, returns the username.\\
                API checks the validity of the input, if the ID is valid, API respond with the username.
            \item get user information\\
                This API recieves the ID of user, returns the user information.\\
                API checks the validity of the input, if the ID is valid, API respond with all user information.
            \item update the user information\\
                This API recieves the ID and information to update, to update user information in the database.\\
                API checks the validity of input, if ID is invalid, or one of all the parameter is empty, the API throws an exception.
            \item course information\\
                This API need the login data of user. Requests backend to respond all courses the user has taken.
            \item course delete\\
                This API need the login data of user. send the id of the course to be deleted, requests backend to delete this course information in user's take table.
                If the user doesn't take this course, this API throws an exception.
        \end{itemize}
\end{itemize}
\subsection{Communications Interfaces}
The data transferring in the Internet is video, and we choose the broadcast protocol is http-flv. This protocol is widely ued in the living industry, with the advantage of low-delay and ease to go through the firewall, it is also friendly to dispatch the rate of flow. \\ 
The session communication data is transferred in form of json. Json is a lightweight data-interchange format. It is easy for humans to read and write. It is easy for machines to parse and generate. Using jwt to authenticate can confirm safety of user session data. JSON Web Token (JWT) is an Internet standard for creating JSON-based access tokens that assert some number of claims. For example, a server could generate a token that has the claim "logged in as admin" and provide that to a client. The client could then use that token to prove that it is logged in as admin. The tokens are signed by one party's private key (usually the server's), so that both parties (the other already being, by some suitable and trustworthy means, in possession of the corresponding public key) are able to verify that the token is legitimate. The tokens are designed to be compact, URL-safe, and usable especially in a web-browser single-sign-on (SSO) context. JWT claims can be typically used to pass identity of authenticated users between an identity provider and a service provider, or any other type of claims as required by business processes. \\
$<$ $>$
%\section{User Interfaces}
%$<$Describe the logical characteristics of each interface between the software 
%product and the users. This may include sample screen images, any GUI standards 
%or product family style guides that are to be followed, screen layout 
%constraints, standard buttons and functions (e.g., help) that will appear on 
%every screen, keyboard shortcuts, error message display standards, and so on.  
%Define the software components for which a user interface is needed. Details of 
%the user interface design should be documented in a separate user interface 
%specification.$>$

\section{Functional Requirements}
$<$ $>$
%\section{Hardware Interfaces}
%$<$Describe the logical and physical characteristics of each interface between 
%the software product and the hardware components of the system. This may include 
%the supported device types, the nature of the data and control interactions 
%between the software and the hardware, and communication protocols to be 
%used.$>$

\section{Behavior Requirements}
$<$ $>$

%\section{Software Interfaces}
%$<$Describe the connections between this product and other specific software 
%components (name and version), including databases, operating systems, tools, 
%libraries, and integrated commercial components. Identify the data items or 
%messages coming into the system and going out and describe the purpose of each.  
%Describe the services needed and the nature of communications. Refer to 
%documents that describe detailed application programming interface protocols.  
%Identify data that will be shared across software components. If the data 
%sharing mechanism must be implemented in a specific way (for example, use of a 
%global data area in a multitasking operating system), specify this as an 
%implementation constraint.$>$

%\section{Communications Interfaces}
%$<$Describe the requirements associated with any communications functions 
%required by this product, including e-mail, web browser, network server 
%communications protocols, electronic forms, and so on. Define any pertinent 
%message formatting. Identify any communication standards that will be used, such 
%as FTP or HTTP. Specify any communication security or encryption issues, data 
%transfer rates, and synchronization mechanisms.$>$


%\chapter{System Features}
%$<$This template illustrates organizing the functional requirements for the 
%product by system features, the major services provided by the product. You may 
%prefer to organize this section by use case, mode of operation, user class, 
%object class, functional hierarchy, or combinations of these, whatever makes the 
%most logical sense for your product.$>$
%
%\section{System Feature 1}
%$<$Don’t really say “System Feature 1.” State the feature name in just a few 
%words.$>$
%
%\subsection{Description and Priority}
%$<$Provide a short description of the feature and indicate whether it is of 
%High, Medium, or Low priority. You could also include specific priority 
%component ratings, such as benefit, penalty, cost, and risk (each rated on a 
%relative scale from a low of 1 to a high of 9).$>$
%
%\subsection{Stimulus/Response Sequences}
%$<$List the sequences of user actions and system responses that stimulate the 
%behavior defined for this feature. These will correspond to the dialog elements 
%associated with use cases.$>$
%
%\subsection{Functional Requirements}
%$<$Itemize the detailed functional requirements associated with this feature.  
%These are the software capabilities that must be present in order for the user 
%to carry out the services provided by the feature, or to execute the use case.  
%Include how the product should respond to anticipated error conditions or 
%invalid inputs. Requirements should be concise, complete, unambiguous, 
%verifiable, and necessary. Use “TBD” as a placeholder to indicate when necessary 
%information is not yet available.$>$
%
%$<$Each requirement should be uniquely identified with a sequence number or a 
%meaningful tag of some kind.$>$
%
%REQ-1:	REQ-2:
%
%\section{System Feature 2 (and so on)}


\chapter{Other Nonfunctional Requirements}
\label{Other Nonfunctional Requirements}
\section{Performance Requirements}
$<$If there are performance requirements for the product under various 
circumstances, state them here and explain their rationale, to help the 
developers understand the intent and make suitable design choices. Specify the 
timing relationships for real time systems. Make such requirements as specific 
as possible. You may need to state performance requirements for individual 
functional requirements or features.$>$

\section{Safety and Security Requirements}
$<$Specify those requirements that are concerned with possible loss, damage, or 
harm that could result from the use of the product. Define any safeguards or 
actions that must be taken, as well as actions that must be prevented. Refer to 
any external policies or regulations that state safety issues that affect the 
product’s design or use. Define any safety certifications that must be 
satisfied.$>$
%\section{Security Requirements}
$<$Specify any requirements regarding security or privacy issues surrounding use 
of the product or protection of the data used or created by the product. Define 
any user identity authentication requirements. Refer to any external policies or 
regulations containing security issues that affect the product. Define any 
security or privacy certifications that must be satisfied.$>$

\section{Software Quality Attributes}
$<$Specify any additional quality characteristics for the product that will be 
important to either the customers or the developers. Some to consider are: 
adaptability, availability, correctness, flexibility, interoperability, 
maintainability, portability, reliability, reusability, robustness, testability, 
and usability. Write these to be specific, quantitative, and verifiable when 
possible. At the least, clarify the relative preferences for various attributes, 
such as ease of use over ease of learning.$>$

%\section{Business Rules}
%$<$List any operating principles about the product, such as which individuals or 
%roles can perform which functions under specific circumstances. These are not 
%functional requirements in themselves, but they may imply certain functional 
%requirements to enforce the rules.$>$


\chapter{Other Requirements}
\label{Other Requirements}
$<$Define any other requirements not covered elsewhere in the SRS. This might 
include database requirements, internationalization requirements, legal 
requirements, reuse objectives for the project, and so on. Add any new sections 
that are pertinent to the project.$>$

\section{Data Dictionary Requirements}
$<$ $>$

%\section{Appendix A: Glossary}
%%see https://en.wikibooks.org/wiki/LaTeX/Glossary
%$<$Define all the terms necessary to properly interpret the SRS, including 
%acronyms and abbreviations. You may wish to build a separate glossary that spans 
%multiple projects or the entire organization, and just include terms specific to 
%a single project in each SRS.$>$

\section{Appendix A: Analysis Models}
$<$Optionally, include any pertinent analysis models, such as data flow 
diagrams, class diagrams, state-transition diagrams, or entity-relationship 
diagrams.$>$

\section{Appendix B: To Be Determined List}
$<$Collect a numbered list of the TBD (to be determined) references that remain 
in the SRS so they can be tracked to closure.$>$

\end{document}
